\section{Characterisation}\label{sec:dc_char}

\subsection{Ideal Transfer}
An \(n\)-bit code \(\langle b_{n-1} \ldots b_1 b_0 \rangle\) has the integer value
\(D = \sum_{k=0}^{n-1} b_k 2^k\), \(0 \leq D \leq 2^n - 1\). With full scale
reference \(X_{REF}\):
\begin{align*}
  X_{LSB} &= \frac{X_{REF}}{2^n} \\
  \text{DAC:} \hspace{0.5cm} X_{OUT}(D) &= D \cdot X_{LSB} = X_{REF} \sum_{k=0}^{n-1} b_k 2^{k-n} \\
  \text{ADC:} \hspace{0.5cm} D(X_{IN}) &= \left\lfloor \frac{X_{IN}}{X_{LSB}} + \frac{1}{2} \right\rfloor \hspace{1cm} [\text{transitions at } (k - \tfrac{1}{2}) X_{LSB}]
\end{align*}

\subsection{Quantization Noise}
For a busy input the ADC error \(e = D X_{LSB} - X_{IN}\) is uniformly
distributed over \(\left[-\frac{X_{LSB}}{2}, \frac{X_{LSB}}{2}\right]\) with
density \(\frac{1}{X_{LSB}}\):
\begin{align*}
  \overline{e^2} = \int_{-X_{LSB}/2}^{X_{LSB}/2} e^2 \frac{1}{X_{LSB}} de
  = \frac{1}{X_{LSB}} \left[\frac{e^3}{3}\right]_{-X_{LSB}/2}^{X_{LSB}/2}
  = \frac{1}{X_{LSB}} \cdot \frac{2}{3} \cdot \frac{X_{LSB}^3}{8} = \frac{X_{LSB}^2}{12}
\end{align*}
A full scale sine \(x(t) = \frac{X_{REF}}{2}\sin(\omega t)\) has power
\(\frac{X_{REF}^2}{8}\):
\begin{align*}
  SQNR &= \frac{X_{REF}^2 / 8}{X_{LSB}^2 / 12} = \frac{12}{8} \cdot 2^{2n} = 1.5 \cdot 2^{2n} \\
  SQNR_{dB} &= 10\log_{10}(1.5) + 2n \cdot 10\log_{10}(2) = 1.76 + 6.02\, n \hspace{1cm} \text{[dB]}
\end{align*}
The noise is spread uniformly over \(0 \ldots \frac{f_s}{2}\). Keeping only a
signal band \(f_B\) (oversampling ratio \(OSR = \frac{f_s}{2 f_B}\)) keeps
\(\frac{1}{OSR}\) of it:
\[
  SQNR_{dB} = 6.02\, n + 1.76 + 10\log_{10}(OSR)
\]
(each doubling of \(f_s\) buys \(3\,\)dB, half a bit).

\subsection{Static Errors}
With \(X_k\) the actual output for code \(k\) (DAC) or the actual transition
into code \(k\) (ADC), after removing offset and gain (end-point line):
\begin{align*}
  DNL_k &= \frac{X_{k+1} - X_k}{X_{LSB}} - 1 \hspace{1cm} \text{[step size error]} \\
  INL_k &= \frac{X_k - k X_{LSB}}{X_{LSB}} = \sum_{j=0}^{k-1} DNL_j \hspace{1cm} \text{[deviation from the line]}
\end{align*}
\(DNL > -1\) everywhere guarantees a monotonic DAC / no missing codes ADC.
An INL bounded by \(\pm\frac{1}{2}\) LSB is the usual linearity specification.

\subsection{Dynamic Errors}
\begin{itemize}
  \item \textbf{SINAD} (signal to noise and distortion) and
    \textbf{ENOB}: \(ENOB = \frac{SINAD_{dB} - 1.76}{6.02}\).
  \item \textbf{Jitter}: a sampling instant error \(t_j\) (rms) on
    \(x = A\sin(2\pi f_{in} t)\) gives \(e = \frac{dx}{dt} t_j\), with rms slope
    \(\frac{2\pi f_{in} A}{\sqrt{2}}\):
    \begin{align*}
      \overline{e^2} &= \frac{(2\pi f_{in} A)^2}{2} t_j^2 \hspace{1cm}
      SNR_{jitter} = \frac{A^2/2}{\overline{e^2}} = \frac{1}{(2\pi f_{in} t_j)^2} \\
      SNR_{jitter,dB} &= -20\log_{10}(2\pi f_{in} t_j)
    \end{align*}
    e.g. \(f_{in} = 100\,\)MHz, \(t_j = 1\,\)ps: \(64\,\)dB, about 10.4 bits.
  \item \textbf{SFDR}: ratio of the signal to the largest spur
    (Chapter~\ref{ch:analysis}, Q8: harmonics of a non-linear transfer).
\end{itemize}

\subsection{Mismatch}
Matched elements of value \(X\) have independent random errors with
relative standard deviation \(S_X = SD\left(\frac{dX}{X}\right)\) (Pelgrom:
\(S_X \propto \frac{1}{\sqrt{\text{area}}}\), Section~\ref{sec:theory_general_probability}).
A sum of \(m\) unit elements has standard deviation \(\sqrt{m}\, S_X X\) since
variances add.
